\documentclass[10pt,a4paper]{amsart}
\usepackage{amsmath,amssymb,amsthm,graphicx}
% Retained minimal preamble for the Stacks Project chapter flat.tex (More on Flatness).
% Verbatim excerpt of the chapter preamble: Xy-pic setup (native preamble lines 14-18) and the native
% theorem declarations with the mathematical macros (native preamble lines 154-217). Copied mechanically;
% nothing added, removed or redefined. The wrapper class and the amsmath/amssymb/amsthm/graphicx base
% packages are supplied by the assembler, so the class/externaldocument part of the native preamble is omitted.
% For commutative diagrams we use Xy-pic
\usepackage[all]{xy}

% We use 2cell for 2-commutative diagrams.
\xyoption{2cell}
% Theorem environments.
%
\theoremstyle{plain}
\newtheorem{theorem}[subsection]{Theorem}
\newtheorem{proposition}[subsection]{Proposition}
\newtheorem{lemma}[subsection]{Lemma}

\theoremstyle{definition}
\newtheorem{definition}[subsection]{Definition}
\newtheorem{example}[subsection]{Example}
\newtheorem{exercise}[subsection]{Exercise}
\newtheorem{situation}[subsection]{Situation}

\theoremstyle{remark}
\newtheorem{remark}[subsection]{Remark}
\newtheorem{remarks}[subsection]{Remarks}

\numberwithin{equation}{subsection}

% Macros
%
\def\lim{\mathop{\mathrm{lim}}\nolimits}
\def\colim{\mathop{\mathrm{colim}}\nolimits}
\def\Spec{\mathop{\mathrm{Spec}}}
\def\Hom{\mathop{\mathrm{Hom}}\nolimits}
\def\Ext{\mathop{\mathrm{Ext}}\nolimits}
\def\SheafHom{\mathop{\mathcal{H}\!\mathit{om}}\nolimits}
\def\SheafExt{\mathop{\mathcal{E}\!\mathit{xt}}\nolimits}
\def\Sch{\mathit{Sch}}
\def\Mor{\mathop{\mathrm{Mor}}\nolimits}
\def\Ob{\mathop{\mathrm{Ob}}\nolimits}
\def\Sh{\mathop{\mathit{Sh}}\nolimits}
\def\NL{\mathop{N\!L}\nolimits}
\def\CH{\mathop{\mathrm{CH}}\nolimits}
\def\proetale{{pro\text{-}\acute{e}tale}}
\def\etale{{\acute{e}tale}}
\def\QCoh{\mathit{QCoh}}
\def\Ker{\mathop{\mathrm{Ker}}}
\def\Im{\mathop{\mathrm{Im}}}
\def\Coker{\mathop{\mathrm{Coker}}}
\def\Coim{\mathop{\mathrm{Coim}}}

% Boxtimes
%
\DeclareMathSymbol{\boxtimes}{\mathbin}{AMSa}{"02}

%
% Macros for moduli stacks/spaces
%
\def\QCohstack{\mathcal{QC}\!\mathit{oh}}
\def\Cohstack{\mathcal{C}\!\mathit{oh}}
\def\Spacesstack{\mathcal{S}\!\mathit{paces}}
\def\Quotfunctor{\mathrm{Quot}}
\def\Hilbfunctor{\mathrm{Hilb}}
\def\Curvesstack{\mathcal{C}\!\mathit{urves}}
\def\Polarizedstack{\mathcal{P}\!\mathit{olarized}}
\def\Complexesstack{\mathcal{C}\!\mathit{omplexes}}
% \Pic is the operator that assigns to X its picard group, usage \Pic(X)
% \Picardstack_{X/B} denotes the Picard stack of X over B
% \Picardfunctor_{X/B} denotes the Picard functor of X over B
\def\Pic{\mathop{\mathrm{Pic}}\nolimits}
\def\Picardstack{\mathcal{P}\!\mathit{ic}}
\def\Picardfunctor{\mathrm{Pic}}
\def\Deformationcategory{\mathcal{D}\!\mathit{ef}}

% --- appended 2026-10-12 (Chief): the project own definitions of the reference/slogan/history
% display environments, copied mechanically from the shared Stacks preamble lines 8-12
% (\usepackage{verbatim} plus the three \newenvironment lines, which the chapter preamble keeps
% verbatim from the project preamble). Nothing added, removed or redefined: these environments
% carry no mathematical content and the project itself typesets them as comments. Without these
% three lines a counted statement whose own text uses them fails with Environment slogan undefined.
\usepackage{verbatim}
\newenvironment{reference}{\comment}{\endcomment}
\newenvironment{slogan}{\comment}{\endcomment}
\newenvironment{history}{\comment}{\endcomment}
\usepackage{cleveref}
\crefname{originallemma}{original lemma}{original lemmas}
\Crefname{originallemma}{Original lemma}{Original lemmas}
\makeatletter
\newlabel{topology-lemma-constructible-hausdorff-quasi-compact}{{0901}{source}{}{sourcepointer-1}{}}
\newlabel{topology-lemma-constructible-hausdorff-quasi-compact@cref}{{[originallemma][0][]0901}{[1][0][]source}{}{}{}}
\makeatother
\crefname{originallemma}{original lemma}{original lemmas}
\Crefname{originallemma}{Original lemma}{Original lemmas}
\makeatletter
\newlabel{algebra-lemma-spec-spectral}{{090M}{source}{}{sourcepointer-2}{}}
\newlabel{algebra-lemma-spec-spectral@cref}{{[originallemma][0][]090M}{[1][0][]source}{}{}{}}
\makeatother
\title{For a family of finitely generated ideals covering the punctured spectrum of a local domain, the intersection is zero}
\author{The Stacks Project Authors}
\date{stacks-project:flat.tex:local-snapshot-20261011; GNU FDL 1.2 or later}
\begin{document}
\maketitle
\noindent\textit{Editorial note.} All delivered mathematical text is the author's own, copied from the pinned Stacks Project chapter file flat.tex; no statement or proof step was rewritten, reordered or removed.\par\medskip
\noindent\textit{Editorial note.} The preamble is the chapter's own retained minimal preamble (sources/stacks-flat/retained-preamble.tex), a verbatim excerpt of the native Xy-pic setup and of the native theorem declarations and mathematical macros, followed by the project own definitions of the reference, slogan and history display environments (verbatim package plus three newenvironment lines) that the retained file carries at its end and that the project itself typesets as comments. No macro was invented or redefined, and the wrapper class and the amsmath/amssymb/amsthm/graphicx base packages are supplied by the assembler.\par\medskip
\noindent\textit{Editorial note.} The counted statement and proof contain no cross reference, no citation and no included file, so the delivered body is self contained.\par\medskip
\noindent\textit{Editorial note.} Canonical location of the counted proposition: Stacks tag 0ATA (sources/stacks-flat/pinned-original-source.tex line 7344).\par\medskip
\noindent\textit{Allowed standard prerequisites (editorial).}\begin{itemize}
\item The constructible topology on the spectrum of a ring is quasi-compact and Hausdorff. Topology, Lemma labelled topology-lemma-constructible-hausdorff-quasi-compact (Stacks tag 0901) with Algebra, Lemma labelled algebra-lemma-spec-spectral (Stacks tag 090M), invoked at proof lines 7358-7361. Supplies that V(f) is quasi-compact in the constructible topology, which turns the covering of V(f) by the V(I) into a finite subcover.
\item Finitely generated ideals and the relation between vanishing loci and powers of an ideal. Algebra, the elementary facts that V(f) contained in V(I) for finitely generated I gives I\textasciicircum{}m contained in (f) for some m. Supplies the power containment used at proof lines 7364-7365; the reduction of the finite cover to a single ideal is inside the counted span.
\end{itemize}\medskip
\ifdefined\hypertarget\hypertarget{sourcepointer-1}{}\fi
\noindent\textit{Source reference.} Original lemma 0901: sources/stacks-chow/tags.txt, source lines 11394--11394. Cross-chapter standard prerequisite located in the project tag list staged in the local snapshot; invoked at proof line 7359 of the counted proof. It is retained only as a navigation pointer; every load-bearing step is inside the delivered span.\par\medskip
\ifdefined\hypertarget\hypertarget{sourcepointer-2}{}\fi
\noindent\textit{Source reference.} Original lemma 090M: sources/stacks-chow/tags.txt, source lines 11415--11415. Cross-chapter standard prerequisite located in the project tag list staged in the local snapshot; invoked at proof line 7361 of the counted proof. It is retained only as a navigation pointer; every load-bearing step is inside the delivered span.\par\medskip
\begin{lemma}
\label{lemma-big-intersection-is-zero}
Let $A$ be a local domain which is not a field.
Let $S$ be a set of finitely generated ideals of $A$.
Assume that $S$ is closed under products and such that
$\bigcup_{I \in S} V(I)$ is the complement of the generic point of $\Spec(A)$.
Then $\bigcap_{I \in S} I = (0)$.
\end{lemma}

\begin{proof}
Since $\mathfrak m_A \subset A$ is not the generic point of $\Spec(A)$
we see that $I \subset \mathfrak m_A$ for at least one $I \in S$.
Hence $\bigcap_{I \in S} I \subset \mathfrak m_A$.
Let $f \in \mathfrak m_A$ be nonzero. Then
$V(f) \subset \bigcup_{I \in S} V(I)$.
Since the constructible topology on $V(f)$ is quasi-compact
(Topology, Lemma \ref{topology-lemma-constructible-hausdorff-quasi-compact}
and
Algebra, Lemma \ref{algebra-lemma-spec-spectral})
we find that $V(f) \subset V(I_1) \cup \ldots \cup V(I_n)$
for some $I_j \in S$. Because $I_1 \ldots I_n \in S$ we see that
$V(f) \subset V(I)$ for some $I$. As $I$ is finitely generated
this implies that $I^m \subset (f)$ for some $m$ and since
$S$ is closed under products we see that $I \subset (f^2)$ for
some $I \in S$. Then it is not possible to have $f \in I$.
\end{proof}

\end{document}
